% Bilingual methods paper for NOETHER-ISV-v021. UTF-8/LF; build with XeLaTeX.
% The complete Interslavic text is primary; the complete English counterpart follows it.
\documentclass[11pt]{article}
\usepackage{fontspec}
\IfFontExistsTF{Times New Roman}{\setmainfont{Times New Roman}}{\setmainfont{TeX Gyre Termes}}
\IfFontExistsTF{Arial}{\setsansfont{Arial}}{\setsansfont{TeX Gyre Heros}}
\IfFontExistsTF{Consolas}{\setmonofont{Consolas}}{\setmonofont{Latin Modern Mono}}
\usepackage[a4paper,margin=2.25cm]{geometry}
\usepackage{microtype}
\usepackage{amsmath,amssymb}
\usepackage{booktabs,tabularx,array}
\usepackage{enumitem}
\usepackage{xcolor}
\usepackage{fancyhdr}
\usepackage{xurl}
\usepackage[hidelinks,unicode]{hyperref}
\usepackage{bookmark}
\definecolor{NoetherBlue}{HTML}{1C476B}
\hypersetup{
  pdftitle={Od jezyčnogo kanona do povtorimoj publikacije: mašinno-pomagana medžuslovjanska redakcija matematičnyh děl Emmy Noether / From a language canon to a reproducible publication: machine-assisted Interslavic editing of Emmy Noether's mathematical works},
  pdfauthor={Machine-assisted Interslavic editorial project},
  pdfsubject={Bilingual methods, linguistic normalization, deterministic script derivation, validation, and reproducibility for NOETHER-ISV-v021},
  pdfkeywords={Interslavic, Medžuslovjansky, Emmy Noether, AI-assisted editing, linguistic normalization, reproducibility, Cyrillic projection}
}
\setlength{\parindent}{1.15em}
\setlength{\parskip}{0.25em}
\setlist{nosep,leftmargin=1.6em}
\emergencystretch=4em
\raggedbottom
\pagestyle{fancy}
\fancyhf{}
\fancyhead[L]{\small NOETHER-ISV-v021 metodologija / methods}
\fancyhead[R]{\small Medžuslovjansky / Interslavic}
\fancyfoot[C]{\thepage}
\newcommand{\code}[1]{\texttt{#1}}
\newcommand{\hash}[1]{\begingroup\ttfamily\footnotesize\path{#1}\endgroup}
\newcommand{\doiurl}[1]{\href{https://doi.org/#1}{#1}}
\newcommand{\repo}{\href{https://github.com/KokunoYumeto/emmy-noether-isv}{KokunoYumeto/emmy-noether-isv}}
\newcolumntype{Y}{>{\raggedright\arraybackslash}X}

\title{\sffamily\bfseries\color{NoetherBlue}
Od jezyčnogo kanona do povtorimoj publikacije:\\
metodologija mašinno-pomaganoj medžuslovjanskoj redakcije\\[0.65em]
\large From a language canon to a reproducible publication:\\
machine-assisted Interslavic editing of Emmy Noether's mathematical works}
\author{Machine-assisted Interslavic editorial project}
\date{NOETHER-ISV-v021 - 24 avgusta / August 2026\\
Version DOI: \doiurl{10.5281/zenodo.22077468}\\
Concept DOI: \doiurl{10.5281/zenodo.21926382}}

\begin{document}
\maketitle

\begin{abstract}
\noindent\textbf{Medžuslovjansky.}
Tuta statja dokumentuje tvorjenje polnogo medžuslovjanskogo izdanja matematičnyh děl Emmy Noether v redakcijnoj latinici i deterministično izvedenoj kirilici. Cělj jest praktično čitanje govoriteljami slovjanskyh jezykov, uključajuči čitateljev bez anglijskogo jezyka, a ne sama demonstracija jezyčnogo eksperimenta. Izdanje v021 vnosit 4\,053 točno lokalizovanyh popravok v latiničny kanon kroz 12 dodavajemyh rěšenij. Prvyh 656 slučajev popravljajut morfologiju, terminologiju i leksičnu doslědnost; 1\,599 slučajev normalizujut opcionalne soglasniky \code{ć/đ} do osnovnogo zapisa \code{č/dž}; 1\,798 slučajev normalizujut šest drugih opcionalnyh bukv, pri čem četyri inojezyčne \code{è} ostavajut ohranjene. Zamrznuty rezultat imaje nulu nerimskyh leksičnyh blokatorov. Projektor v8 autentifikuje zamrznutje, rebazovany manifest od 5\,290 zapisov, 372 nasledžene efekty i 722 rimsky identitetne lokatory. Nezavisny izvod povtorjaje četyri kirilične fajly bajt-po-bajtu; 2\,049 izměnjenyh vhodnyh blokov zakonno objašnjavajut vse izhodne razlike, a 2\,053 neizměnjene kontrole ostavajut identične. Dva polna čitateljnika imajut 565 i 588 stranic; relizna kontrola imaje 38 uspěšnyh prověrok i nulu pogrěšek. ``Nula blokatorov'' označaje samo završeny deklarovany auditny frontier, ne rodno-govoriteljsku, obćinsku ili naučnu certifikaciju. Vse iztočniki, rěšenja, programy, kontrolne sumy, ograničenja i javne adresy sut ostavljene tako, da druga sesija ili recenzent može povtoriti, kritikovati ili monotono prodolžiti izdanje.

\medskip
\noindent\textbf{English.}
This paper documents the production of a complete Interslavic edition of Emmy Noether's mathematical works in editorial Latin script and deterministically derived Cyrillic. The purpose is practical reading by speakers of Slavic languages, including readers without English, rather than a display of linguistic experimentation. Release v021 makes 4,053 exactly localized corrections to the Latin canon through 12 append-only decisions. The first 656 sites repair morphology, terminology, and lexical consistency; 1,599 sites normalize optional consonants \code{ć/đ} to the basic surface \code{č/dž}; and 1,798 sites normalize six other optional letters while preserving four foreign \code{è} spellings. The frozen result has zero non-Roman lexical projection blockers. Projector v8 authenticates that freeze, a rebased 5,290-record manifest, 372 inherited effects, and 722 Roman-identity loci. An independent derivation reproduces the four Cyrillic files byte for byte; 2,049 changed input blocks license every output delta, while 2,053 unchanged controls remain identical. The complete readers contain 565 and 588 pages; release QA reports 38 passing checks and zero errors. ``Zero blockers'' denotes only the completed declared audit frontier, not native-speaker, community, or scholarly certification. Sources, decisions, software, hashes, limitations, and public access points are retained so another session or reviewer can reproduce, challenge, or monotonically extend the edition.
\end{abstract}

\part*{Čest I - Polny medžuslovjansky tekst}
\addcontentsline{toc}{part}{Čest I - Polny medžuslovjansky tekst}

\section{Predmet, cělj i čitatelji}

Predmet jest jedno polno medžuslovjansko jezyčno izdanje ograničenogo korpusa Emmy Noether: numerovane Statje 1--43, lekcijno dělo 44, Statja 45 i bibliografija. To ne jest antologija priměrov i ne jest zbirka nevezanyh transkripcij. Javny čitateljnik dava dva pisemne vida jedinogo jezyčnogo objekta:

\begin{itemize}
\item \code{isv-Latn} jest kanoničny redakcijny jezyčny tekst;
\item \code{isv-Cyrl} jest deterministična projekcija togo teksta, a ne nezavisny svědok prěvoda.
\end{itemize}

Toto razděljenje jest metodologično važno. Izvedena azbuka ne može sama potvrditi pravilnost prěvoda, iz ktorego jest stvorjena. Jezyčne popravky vhodet jedino v latinny kanon; kirilica mora byti ponovno izvedena, testovana, sgrađena i vizualno prověrjena po zamrznutju vhodnyh bajtov. V021 jest osobito jasny slučaj: latiničny kanon dostavaje materialne jezyčne popravky i osnovnu čitateljsku ortografiju, potom kirilica se ponovno izvodi iz cělogo točnogo vektora, a ne popravljaje ručno kako konkurentny svědok.

Praktična heuristika jezyčnogo napredka ne jest pytanje ``koliko tekst izgleda slovjansky?'' Uporabjajemo sojedinjenje četyrěh pytanj:

\begin{enumerate}
\item Može li govoritelj različnogo slovjanskogo jezyka pravděpodobno razobrati větu bez anglijskogo teksta?
\item Jest li izbrana forma podprta deklarovanym medžuslovjanskim slovnikom ili gramatikoju, albo jasno zapisanym domovym pravilom projekta?
\item Ohranjaje li forma točny matematičny predmet, kvantifikaciju i odnos argumentov německogo iztočnika?
\item Jest li vsaky slučaj izključka, identiteta, citaty ili nejasnosti reprezentovany povtorimoju roloju, vměsto bezglasnogo izostavjenja?
\end{enumerate}

Zato izdanje jest najbolje razuměti kako popravljivo naučno rabotno izdanje: javno, pregledljivo, vezano kontrolnymi sumami i otvorjeno za ulučšenje, ale bez tvrdženja o rodno-govoriteljskoj ili kritično-izdateljskoj certifikaciji.

\section{Arhitektura avtoriteta i kustodije}

Metoda razděljaje iztočny avtoritet, jezyčny kanon, izvedenu azbuku, redakcijne dokazy i javne produkte. Nijedin sloj ne smije samovoljno prěvzeti funkciju drugogo.

\begin{table}[ht]
\centering
\caption{Role avtoriteta v NOETHER-ISV-v021.}\label{tab:isv-authority}
\begin{tabularx}{\textwidth}{@{}p{0.21\textwidth}YY@{}}
\toprule
Rola & Upravljajuči objekt & Poslědstvo\\
\midrule
Německy sravnjivateljny svědok & Objekt \code{NOETH-DE-ED-0015}, upotrěbjen pri redakciji v021 & Upravjaje iztočnym smyslom, strukturoju i matematičnym sravnjenjem; medžuslovjanska linija go ne rediguje niti tvrdi, že on jest najnovějši německy reliz.\\
Latiničny jezyčny kanon & Četyri \code{source\_latin/*.tex} fajly pri glave \code{ISV021-EDIT-0012} & Jedino město jezyčnyh redakcij; v021 priměnilo 4\,053 lokalizovane popravky.\\
Manifest rolij & \hash{ISV021-REBASED-PROJECTION-ROLE-MANIFEST.json} & Veže 5\,290 zapisov, uključajuči 1\,980 zaščitnyh granic i 722 rimsky lokatory, k točnym iztočnym mjestam.\\
Kirilična površina & Projektor v8 nad točnym latinnym vektorom i manifestom & Povtorimo tvorjeny izhod; ne može nadglasovati latinny kanon niti nezavisno potvrditi prěvod.\\
Relizne produkty & Dva čitateljnika, vezany čitateljnik, dva ZIP-arhiva i SHA-256 manifest & Člověku dostupno razprostranjenje; produkt ne može redefinovati kanon, iz ktorego byl stvorjen.\\
\bottomrule
\end{tabularx}
\end{table}

Německy objekt naslědženy iz v019 imaje SHA-256
\hash{51A25101C04877AE740989E72B2AD65A7A7E65B081077C4A518BF1737AD5B907}.
Ta identiteta označaje kontrolny iztočnik prěma kojemu latinny kanon byl zamrznuty; ona ne jest tvrdženje, že projekt publikoval novy kritičny německy tekst. Globalna mnogojezyčna zbirka jest povezana odnosom \code{IsPartOf} s \doiurl{10.5281/zenodo.20412587}, a identificirany německy objekt odnosom \code{IsTranslationOf}.

\section{Mašinno čitljivo odkryvanje kanona}

Povtorimost začinaje s odkryvanjem. Buduća sesija ne smije ugadyvati ktory iz podobnyh fajlov jest kanon. Zato koren repozitorija imaje \code{CANON\_INDEX.json}, ktory deklaruje:

\begin{itemize}
\item jezyk i BCP~47 označenja pisma;
\item četyri kanonične latinne fajly i četyri izvedene kirilične fajly;
\item točne bajtove i SHA-256 vsakej upravljajučej jedinice;
\item latinnu redakcijnu glavu i identitet žurnala, do ktorogo se može samo dodavati;
\item manifest rolij, projektor, test i prověrčny dokument;
\item německy avtoritet, granice pisanja i granice dozvoljenyh tvrdženj;
\item deterministične komandy sgrađanja i imena javnyh artefaktov.
\end{itemize}

Indeks jest deklarativny, a ne samopotvrđajuči. On govori agentu čto objekty \emph{sut}; izvršive prověrky dokazujut, jest li deklarovana identiteta ješče pravdiva. Tuto razděljenje odklanjaje kružno pravilo ``najnověje izgledajuči fajl jest kanon'' i dava osnovu, k ktoroj druga transkripcija, prěvod ili azbučna projekcija može bezpečno pridati novy sloj.

\section{Čto AI děla i čto ne može rěšiti samo}

Projekt uporablja AI kako redakcijny instrument, ne kako iztočny avtoritet. Mašinny trud obimaje sravnjenje větnyh kontekstov, tvorjenje kandidatnyh inventarov, klasifikaciju pojavjenj, formulovanje alternativ, generovanje strukturnyh testov i tvorjenje javnoj dokumentacije. Deterministične programy potom autentifikuju vhody, priměnjajut samo zamrznute rěšenja i odklanjajut nepredvidženu izměnu.

Metoda ima slědujuči cikl:

\begin{enumerate}
\item zamrznuti iztočny i ciljevy bajtovy vektor;
\item najdti materialne jezyčne ili rolne kandidaty bez avtomatičnogo globalnogo zaměnjenja;
\item pregledati vsako pojavjenje v větnom, bibliografičnom, matematičnom i TeX-kontekstu;
\item zapisati izbranu rolu, dokaz, odklonjene alternativy, neuvěrjenost i poslědstva;
\item napraviti manifest ili redakcijno rěšenje, vezano k točnomu lokatoru i kontrolnoj sumě;
\item izvesti izhod čistym programom i dokazati dopuščeny obseg razlike;
\item sgraditi, iztegnuti tekst, prověriti fonty, navigaciju, strukturu i vizualne stranice;
\item publikovati izvorniki, dokazy i haše zajedno s čitateljnikom, potom anonimno pročitati javne bajty nazad.
\end{enumerate}

AI-prědlog ne stavaje kanonom zato, že zvuči uvěrlivo. On stavaje dopustimym samo kogda jest lokalizovany, podporjen deklarovanym dokazom, ograničeny v obsegu i povtorimo prověrjen. Kde dokaz ne jest dostatočny, pravilen rezultat jest ograničeno tvrdženje ili zapisano nerěšeno pytanje, a ne izmišljena certifikacija.

\section{Korekcija latinnogo kanona i rebase projekcije v021}

V021 začinaje od funkcionalnogo v020 projektora i cělo pregledanogo manifesta rolij, ale ponovno otvara latiničny jezyčny kanon. Kandidaty sut razděljene na morfologiju, smysl, terminologiju, leksičnu doslědnost i jasno deklarovanu standardnu osnovnu ortografiju. Nijedna globalna zaměna ne jest dovoljena bez pojavjenjskogo manifesta, ochrany TeX-zon, točnogo obrata i nezavisnogo prověrenja naslědnika.

\begin{table}[ht]
\centering
\caption{V021 latinična korekcija i projekcijny manifest.}\label{tab:isv-review}
\begin{tabular}{@{}lrr@{}}
\toprule
Sloj & Pojavjenja & Status\\
\midrule
Morfologija, terminologija i leksična doslědnost & 656 & priměnjeno\\
Opcionalne \code{ć/đ} do osnovnogo \code{č/dž} & 1\,599 & priměnjeno\\
Šest drugih opcionalnyh bukv do osnovnogo zapisa & 1\,798 & priměnjeno\\
Rimsky identitetne lokatory & 722 & ohranjeno\\
\midrule
Vse latinične popravky & 4\,053 & 12 rěšenij\\
Vse zapisy rebazovanogo manifesta & 5\,290 & nula blokatorov\\
\bottomrule
\end{tabular}
\end{table}

Poslě deset leksično-morfologičnyh rěšenij slijeduje osnovna čitateljska površina. Opcionalne \code{ć} i \code{đ} sut zaměnjene s \code{č} i \code{dž} v 1\,599 vidimyh slučajev. Potom \code{ė}, \code{ę}, \code{ľ}, \code{ś}, \code{ť} i \code{ų} sut normalizovane v 1\,798 slučajev, ale četyri dokazano inojezyčne \code{è} v imenah i citatah ostavajut netknute. Tuto jest točny allowlist, ne generično odstranjenje diakritiky.

Rebase ohranjaje všech 4\,568 naslědženyh zapisov jedan-k-jednom, v tom 1\,980 ohranjajučih zapisov bajt-po-bajtu. K nim dodava 722 točne rimsky lokatory. V8 ponovno dokazuje svih 372 naslědženyh efektov: 346 ohranjajučih i 26 projicirajučih. Sravnjenje v020--v021 nalazi 2\,049 izměnjenyh blokov, ktore licencujut vse izhodne razlike, i 2\,053 neizměnjene kontrolne bloky s identičnym izhodom. Nezakonnyh izhodnyh izměn jest nula.

Status ``nula blokatorov'' treba čitati operacionno: deklarovany auditny frontier ne imaje nijedin nerimsky leksičny blok projekcije, a vsaky od 722 rimskyh tokenov imaje točny identitetny lokator. On ne govori, že prirodny jezyk ne imaje polisemije, že vsaky govoritelj izbere tu samu formu, ili že někoj rodny govoritelj pregledal cijely korpus.

\section{Jezyčna normalizacija latinnogo kanona s žurnalom, do ktorogo se može samo dodavati}

Latinna glava naslědžuje 171 materialnyh rěšenj iz v019 i dodava 12 rěšenij v021. Rěšenje jest dopustimo samo ako imenuje prědhodnika, točnu ciljnu poziciju, prědlagany naslědnik, relevantny avtoritet, odklonjene alternativy, motivaciju, neuvěrjenost i poslědnu prověrku. Aplikator autentifikuje vsaky vhod i odklanjaje mnogoznačny cilj. V žurnal se može samo dodavati: pogrěšny ili pozdněje zaměnjeny postup ne izčezaje, ale ostavaje čitljivym s novym poslědnikom.

Glavne jezyčne klasy sut:

\begin{itemize}
\item \textbf{Lažni prijatelji i smyslove razděljenja.} Podobnost s ruskym, ukrajinskym, poljskym, češskym ili južnoslovjanskym oblikom sama ne dostaje. Obstajanje, sučnost, sostav, sodržanje, prědpoloženje, poslědstvo i odgovaranje sut rěšene po větnom i matematičnom smyslu.
\item \textbf{Morfologija i větny položaj.} Soglasovanje, kosy padeži, genitiv množiny, pričestja, glagolne koncovky i poredkove formy sut klasifikovane prěd popravkoju. Globalna zaměna koncovky jest odklonjena, ibo formuly, imena i citaty mogut imitovati običny morfem.
\item \textbf{Vezy i logika dokaza.} Uzročne, zaključne, uslovne, ekvivalentne, protivne i apodozne vezy sut sravnjene s odnosom iztočnyh klauz. Čestost slova ne smije obrnuti implikaciju ili oslabiti kvantifikator.
\item \textbf{Imena, naslovi i citaty.} Polny neizměnjeny identitet i izvornojezyčna bibliografična površina sut ohranjene. Integrirana ili sklanjana forma projiciraje se kogda funkcionuje kako medžuslovjanska morfologija.
\end{itemize}

V021 priměnjaje te rěšenja k latinnomu kanonu, zamrzuje točny naslědnik pri \code{ISV021-EDIT-0012}, potom ponovno veže projekcijne role k novym lokatoram. Projekcija nikogda ne jest uporabljena kako skryto město jezyčnoj redakcije.

\section{Matematična i TeX-věrnost}

Prirodnojezyčna normalizacija jest podredžena matematičnoj identitetě. Zaštićene granice obimajut matematičny režim, imena komand, opcije komand, oznaky, odkazy, ključe citat, puti, komentare i dokazano izvornojezyčne argumenty. Za vsaky latinni/kiriličny par odstranjenje Unicode-bukv i kombinovanyh znakov dava tu samu ne-bukvenu struju. Taka prověrka ohranjaje zagrady, kontrolne slova, operatory, čisla, poravnateljne znaky i matematične symbole.

Ranějši neuspěh pokazuje, začto jezyčny skener ne dostaje. Projektor v5 je pogrěšno čital TeX-jedinicu \code{24em} kako vidimo slovo i stvoril kiriličnu jedinicu, ktora ne jest zakonny TeX. Projektor v6 je zato dobil eksplicitnu gramatiku numeričnyh razměrov, ktore ju v8 naslědžuje i ponovno prověrjaje. Ohranjeno jest 447 raněje izloženih pojavjenj: 219 \code{em} i 228 \code{pt}; kiriličny ostatok razměrne jedinice jest nula. Neuspěh ostavaje v dokazah, ibo ``čisty'' leksičny rezultat ne dokazuje izvršivu pravilnost iztočnogo programa.

\section{Determinističny izvod v8}

V8 jest čista funkcija četyrěh zamrznutyh latinnyh fajlov, rebazovanogo manifesta rolij i pinovanyh zavisnostij. On:

\begin{enumerate}
\item autentifikuje razměr i SHA-256 vsakoj latiničnoj jedinice;
\item autentifikuje latiničny freeze-zapis i freeze-rěšenje prěd projekcijeju;
\item prověrjaje točnu topologiju iztočnikov i manifesta;
\item razděljuje TeX na vidime i zaštićene role;
\item sjedinjuje pregledane granice imen, citat, 722 rimskyh označenj i razměrov;
\item projiciraje ostale vidime riječi zamrznutym medžuslovjanskim transliteratorom;
\item dokazuje vsaky iz 372 nasledženyh efektov i vsaky od 722 rimskyh lokatorov;
\item licencuje vsaku izhodnu razliku točnym v020--v021 latinnym blokom;
\item piše izhody transakcijno samo poslě vseh prověrok, a zapis potvrđenja děla v odvojenej transakciji.
\end{enumerate}

Hostilny testovi prověrjajut čisty montáž, idempotentny povtor, odklonjenje tretjego stanja, rollback pri naměrnom drugom instalacijskom neuspěhu i odklonjenje kombinovanoj izhodno-potvrđajučej transakcije prě projekcijeju. Tri identične testove prolazet. Nezavisny prověrnik, ktory ne importuje v8, ponovno tvori četyri kirilične izhody i poslě instalacije čita žive bajty nazad. Test, projektor i prověrnik sami sut hašovani, tako že ``isti postup'' znači identičny program, ne samo podobno opisanje.

\section{Sgrađanje i mnogoslojna kontrola}

Vse XeLaTeX-operacije tečut serijno, s dvoma prohodami po vsakom dokumentu, s izključenym shell escape i fiksovanym \code{SOURCE\_DATE\_EPOCH}. Osm redaktabilnyh komponentov davajut dva polna čitateljnika.

\begin{table}[ht]
\centering
\caption{Prověrjene komponentne čitateljniki v021.}\label{tab:isv-metrics}
\begin{tabularx}{\textwidth}{@{}YrrY@{}}
\toprule
Artefakt & Stranice & Bajty & SHA-256\\
\midrule
Latiničny čitateljnik & 565 & 2\,993\,022 & \hash{57FA897FCADA9C9A8BC764BB0DE30724B7434C31B8C8A35B14FB42922DAB898B}\\
Kiriličny čitateljnik & 588 & 3\,300\,230 & \hash{9513C002794173DD6CC44F157489B0B7A94E92E19378605D3E15C5EE27E895D3}\\
\bottomrule
\end{tabularx}
\end{table}

Relizna kontrola imaje 38 uspěšnyh prověrok i nulu pogrěšek. Ona obimaje točno povratno čitanje iztočnikov, TeX i ne-bukvenu strukturu, fatalne žurnale, straničnu topologiju, fonty, \code{pdftotext}, popis pisem i identitet teksta medžu komponentami i kumulativnymi čitateljnikami. Vezany javny čitateljnik mora imati 6 uvodnyh/razdělnyh stranic plus 565 latiničnyh i 588 kiriličnyh stranic, vkupno 1\,159. Njegov graditelj odklanjaje drugu topologiju.

Strukturny test ne dokazuje krasivy izgled. Zato finalna relizna kontrola takože zahtěvaje serijno renderovanje i vizualny pregled novogo uvoda, vsakej izměnjene stranice i susědnyh kontrolnyh stranic, plus razdělnyh granic dvoh pisem. Tekstny izteg ne može sam odkryti prěkryvanje, obrezanje ili nečitljivu glifnu zaměnu.

\section{Publikacija, dostup i ohranjanje}

Javna površina Zenodo imaje četyri produkta:

\begin{enumerate}
\item \code{00\_NOETHER\_INTERSLAVIC\_COMPLETE\_LINKED\_READER.pdf};
\item \code{01\_NOETHER\_INTERSLAVIC\_EDITABLE\_SOURCES.zip};
\item \code{02\_NOETHER\_INTERSLAVIC\_EVIDENCE\_AND\_PROVENANCE.zip};
\item \code{03\_NOETHER\_INTERSLAVIC\_SHA256\_MANIFEST.txt}.
\end{enumerate}

Prvy fajl jest ljudsky čitateljnik. Iztočny arhiv sodrži latinni kanon, izvedene kirilične izvorniky, projektor, testy, graditelje, metadane i členske manifesty. Dokazny arhiv sodrži redakcijny žurnal, do ktorogo se može samo dodavati, bočne zapisy rěšenj, polny manifest rolij, prověrčne dokumenty, tuta metodologična statja s LaTeX-iztočnikom, QA i zadržane dokazy neuspěhov. Operacijske predpomněti, tajne ključe, surove razgovorne zapisy i nevezane rabotne oblasti ne sut javny dokazny produkt.

Stabilny koncept izdanja jest \doiurl{10.5281/zenodo.21926382}; verzijny DOI v021 jest \doiurl{10.5281/zenodo.22077468}; javny repozitorij jest \repo. Opis v repozitoriju i metadane Zenodo stavjajut najprvo polny medžuslovjansky uvod, potom polnu anglijsku paralelnu verziju. Publikacija završaje se anonimnym povratnym čitanjem vsakego javnogo fajla i sravnjenjem bajtov i SHA-256 s lokalnym zamrznutym manifestom.

\section{Ograničenja, ljudsky pregled i odgovorne tvrdženja}

Izdanje može s uvěrjenostju tvrdžiti točny vhodny vektor, povtorimu redakcijnu historiju, polny pregled projektorskih rolij, determinističny kiriličny izvod, ograničenu izhodnu razliku, sgrađanja bez pogrěšek i prověrjenu PDF-strukturu. Ono ne tvrdži:

\begin{itemize}
\item rodno-govoriteljsku certifikaciju ili soglasje medžu slovjanskymi obćinami;
\item potvrđenje medžuslovjanskoj obćiny ili oficialnogo standarda;
\item nezavisnu naučnu recenziju;
\item novu kritičnu redakciju německyh originalov;
\item že računalna prověrka dokazuje idiomatičnost ili optimalny leksičny izbor;
\item že kirilica jest nezavisny dokaz pravilnosti prěvoda.
\end{itemize}

Člověčsky vlastnik projekta je oprědělil obseg, kustodiju, dostup i publikacijne cěljy. Dostupny dokaz ne opravdava tvrdženje, že člověk ili rodny govoritelj je liniju-po-liniji recenzoval cěly jezyčny korpus. Slova ``završeny auditny frontier'' v v021 označajut završenu deklarovanu pojavjenjsku i projekcijnu kontrolu, izvedenu mašinno-pomaganym trudom i prověrjenu determinističnymi kontrolami. Ona ne značet ljudsku jezyčnu certifikaciju.

Ta ograničenja ne sut razlog skryti dělo. Javny izvorniky, lokatory i kontrolne sumy prětvarjajut nesoglasje v izvršivu popravku: recenzent može imenovati město, prědložiti lučše čitanje, citovati protivny dokaz i napraviti naslědno izdanje, ktore dodava popravku bez izbrisanja prědhodnogo stanja.

\section{Povtorimy protokol za druge jezyky i transkripcije}

Najmenši prěnosimy protokol jest:

\begin{enumerate}
\item oprěděliti jedin iztočny avtoritet i jedin ciljevy jezyčny kanon;
\item pridati eksplicitnu rolu vsakoj izvedenoj azbukě, transkripciji ili formatu;
\item zamrznuti iztočne jedinice s bajtami i hašami prěd redakcijeju;
\item zapisovati materialne rěšenja v žurnal, do ktorogo se može samo dodavati, zajedno s odklonjenymi alternativami;
\item razděliti vidimy jezyk od zaštićenoj sintakse, matematike i identitetov;
\item pregledati kandidaty po pojavjenju i vezati klasifikacije k stabilnym lokatoram;
\item prověriti smysl, jezyk, matematiku, sintaksu, sgrađanje, izteg, fonty i renderovanje kako različne vorota;
\item publikovati ljudsky čitateljnik, redaktabilne izvorniky, dokazy i vrhnji haš-manifest;
\item izložiti mašinno čitljiv indeks kanona i deterministične vhody;
\item javno priznati odsutny pregled i ostalu neuvěrjenost;
\item izdavati neměnljive verzije i anonimno prověriti javne bajty.
\end{enumerate}

Glavna lekcija jest prosta: AI-stvorjeno jezyčno izdanje stavaje naučno uporabljivo togda, kogda njegovi izbori ne ostavajut v razgovoru modela. Kanon, role, dokazy, neuspěhi, komandy sgrađanja i javne identitety morajut živěti nezavisno od sesije, ktora jih je stvorila.

\section*{Citovanje, dostupnost i ponovno uporabjenje}

Prědlagano citovanje: \emph{Machine-assisted Interslavic editorial project. Od jezyčnogo kanona do povtorimoj publikacije: metodologija mašinno-pomaganoj medžuslovjanskoj redakcije / From a language canon to a reproducible publication: machine-assisted Interslavic editing of Emmy Noether's mathematical works. NOETHER-ISV-v021.} \doiurl{10.5281/zenodo.22077468} (2026).

Tuta metodologična statja i druge projektom stvorjene redakcijne materialy sut ponudžene pod CC0 v měrě, v ktorej prava obstajajut. Iztočne děla i materialy tretjih stran zadrživajut svoj vlastny pravny status. LaTeX-iztočnik jest vključen, da tekst može byti popravljany, razširjany ili ponovno publikovany s točnym avtorstvom, verzijnoju identitetoju i odkrytym rodoslovjem dokazov.

\clearpage
\part*{Part II - Complete English counterpart}
\addcontentsline{toc}{part}{Part II - Complete English counterpart}

\section{Object, purpose, and readership}

The object is one complete Interslavic language edition of a bounded Emmy Noether corpus: numbered Papers 1--43, lecture work 44, Paper 45, and the bibliography. It is not an anthology of samples and not a collection of unrelated transcriptions. The public reader provides two writing surfaces of one language object:

\begin{itemize}
\item \code{isv-Latn} is the canonical editorial language text;
\item \code{isv-Cyrl} is a deterministic projection of that text, not an independent translation witness.
\end{itemize}

This distinction is methodologically important. A derived script cannot itself confirm the translation from which it was created. Linguistic corrections enter only the Latin canon; Cyrillic must be regenerated, tested, built, and visually reviewed after the input bytes are frozen. V021 is an especially clear case: the Latin canon receives material language corrections and a basic reader-facing orthography, after which Cyrillic is regenerated from the complete exact vector instead of being hand-corrected as a competing witness.

The practical heuristic for linguistic progress is not ``how Slavic-looking is the prose?'' It is the conjunction of four questions:

\begin{enumerate}
\item Can speakers of different Slavic languages plausibly parse the sentence without English?
\item Is the selected form supported by the declared Interslavic lexicon or grammar, or by a clearly recorded project-house rule?
\item Does the form preserve the exact mathematical referent, quantification, and argument relation of the German source?
\item Is every exception, identity, quotation, or ambiguity represented by a reproducible role rather than silently omitted?
\end{enumerate}

The edition is therefore best understood as a corrigible scholarly working edition: public, inspectable, hash-bound, and open to improvement, but not represented as native-speaker-certified or as a critical edition.

\section{Authority and custody architecture}

The method separates source authority, language canon, derived script, editorial evidence, and public products. No layer may silently assume the function of another.

\begin{table}[ht]
\centering
\caption{Authority roles in NOETHER-ISV-v021.}\label{tab:en-authority}
\begin{tabularx}{\textwidth}{@{}p{0.21\textwidth}YY@{}}
\toprule
Role & Controlling object & Consequence\\
\midrule
German comparison witness & Object \code{NOETH-DE-ED-0015}, used during v021 editing & Controls source sense, structure, and mathematical comparison; the Interslavic lane neither edits it nor claims it is the newest German release.\\
Latin language canon & Four \code{source\_latin/*.tex} files at head \code{ISV021-EDIT-0012} & The only surface for linguistic edits; v021 applies 4,053 localized corrections.\\
Role manifest & \hash{ISV021-REBASED-PROJECTION-ROLE-MANIFEST.json} & Binds 5,290 records, including 1,980 protective boundaries and 722 Roman loci, to exact source locations.\\
Cyrillic surface & Projector v8 over the exact Latin vector and manifest & Reproducibly generated output; cannot overrule the Latin canon or independently validate the translation.\\
Release products & Two readers, linked reader, two ZIP archives, and SHA-256 manifest & Human-facing dissemination; products cannot redefine the canon that produced them.\\
\bottomrule
\end{tabularx}
\end{table}

The German object inherited from v019 has SHA-256
\hash{51A25101C04877AE740989E72B2AD65A7A7E65B081077C4A518BF1737AD5B907}.
This identity marks the control source against which the Latin canon was frozen; it is not a claim that this project published a new critical German text. The global multilingual collection is related by \code{IsPartOf} to \doiurl{10.5281/zenodo.20412587}, and the identified German object by the semantic relation \code{IsTranslationOf}.

\section{Machine-readable canon discovery}

Repeatability begins with discovery. A future session must not infer which of several similar files is canonical. The repository root therefore exposes \code{CANON\_INDEX.json}, which declares:

\begin{itemize}
\item language and BCP~47 script tags;
\item the four canonical Latin files and four derived Cyrillic files;
\item exact byte counts and SHA-256 values for every controlling unit;
\item the Latin editorial head and append-only ledger identity;
\item the role manifest, projector, test, and validation receipt;
\item German authority, write boundaries, and claim boundaries;
\item deterministic build commands and public artifact names.
\end{itemize}

The index is declarative, not self-authenticating. It tells an agent what the objects \emph{are}; executable checks prove whether the declared identities still hold. This separation rejects the circular rule ``the newest-looking file is the canon'' and provides a base to which another transcription, translation, or script projection can safely add a new layer.

\section{What AI does and what it cannot decide alone}

The project uses AI as an editorial instrument, not as source authority. Machine work includes sentence-level comparison, candidate inventory construction, occurrence classification, formulation of alternatives, generation of structural tests, and production of public documentation. Deterministic programs then authenticate inputs, apply only frozen decisions, and reject unintended changes.

The method follows this cycle:

\begin{enumerate}
\item freeze source and target byte vectors;
\item locate material linguistic or role candidates without automatic global replacement;
\item review every occurrence in sentence, bibliography, mathematics, and TeX context;
\item record the chosen role, evidence, rejected alternatives, uncertainty, and consequences;
\item produce a manifest or editorial decision bound to an exact locator and hash;
\item derive output with a pure program and prove the allowed delta;
\item build, extract text, inspect fonts, navigation, structure, and rendered pages;
\item publish sources, evidence, and hashes with the reader, then anonymously read the public bytes back.
\end{enumerate}

An AI proposal does not become canonical because it sounds plausible. It becomes admissible only when localized, supported by declared evidence, bounded in scope, and reproducibly verified. Where evidence is insufficient, the proper output is a narrow claim or a recorded open question, not invented certification.

\section{V021 Latin-canon correction and projection rebase}

V021 starts from the functional v020 projector and its fully reviewed role manifest, but reopens the Latin language canon. Candidates are separated into morphology, sense, terminology, lexical consistency, and a deliberately declared standard basic orthography. No global replacement is admissible without an occurrence manifest, protected TeX zones, exact reversal, and independent successor verification.

\begin{table}[ht]
\centering
\caption{V021 Latin correction and projection manifest.}\label{tab:en-review}
\begin{tabular}{@{}lrr@{}}
\toprule
Layer & Occurrences & Status\\
\midrule
Morphology, terminology, and lexical consistency & 656 & applied\\
Optional \code{ć/đ} to basic \code{č/dž} & 1,599 & applied\\
Six other optional letters to the basic surface & 1,798 & applied\\
Roman-identity loci & 722 & preserved\\
\midrule
All Latin corrections & 4,053 & 12 decisions\\
All rebased-manifest records & 5,290 & zero blockers\\
\bottomrule
\end{tabular}
\end{table}

After ten lexical-morphological decisions, the project fixes a basic reader-facing surface. Optional \code{ć} and \code{đ} become \code{č} and \code{dž} at 1,599 visible sites. Then \code{ė}, \code{ę}, \code{ľ}, \code{ś}, \code{ť}, and \code{ų} are normalized at 1,798 sites, while four demonstrably foreign \code{è} spellings in names and quotations remain untouched. This is an exact allowlist, not generic diacritic stripping.

The rebase preserves all 4,568 inherited records one to one, including 1,980 protective records byte for byte, and adds 722 exact Roman loci. V8 re-proves all 372 inherited effects: 346 preserving and 26 projecting. The v020--v021 alignment finds 2,049 changed blocks that license every output difference and 2,053 unchanged control blocks with identical output. Unlicensed output changes are zero.

The status ``zero blockers'' must be read operationally: the declared audit frontier contains no non-Roman lexical projection blocker, and each of the 722 Roman tokens has an exact identity locator. It does not assert that natural language has no polysemy, that every speaker would choose the same form, or that a native speaker reviewed the complete corpus.

\section{Append-only linguistic normalization of the Latin canon}

The Latin head inherits 171 material decisions from v019 and adds 12 v021 decisions. A decision is admissible only when it identifies a predecessor, exact target locus, proposed successor, relevant authority, rejected alternatives, motivation, uncertainty, and post-application validation. An applicator authenticates every input and rejects ambiguous targets. The ledger is append-only: a failed or superseded approach remains readable with its successor instead of disappearing.

The main linguistic classes are:

\begin{itemize}
\item \textbf{False friends and semantic partitions.} Resemblance to Russian, Ukrainian, Polish, Czech, or South Slavic forms is insufficient. Existence, essence, composition, containment, assumption, consequence, and correspondence are adjudicated by sentence and mathematical sense.
\item \textbf{Morphology and sentence role.} Agreement, oblique cases, genitive plurals, participles, verb endings, and ordinal forms are classified before correction. Global suffix replacement is rejected because formulas, names, and quotations can imitate an ordinary morpheme.
\item \textbf{Connectives and proof logic.} Causal, inferential, conditional, equivalence, contrastive, and apodosis functions are compared with the relation between source clauses. Word frequency must not reverse implication or weaken a quantifier.
\item \textbf{Names, titles, and quotations.} A complete uninflected identity and an original-language bibliographic surface are preserved. An integrated or inflected form is projected when it functions as Interslavic morphology.
\end{itemize}

V021 applies those decisions to the Latin canon, freezes the exact successor at \code{ISV021-EDIT-0012}, and then rebinds projection roles to the new locators. Projection is never used as a hidden surface for language editing.

\section{Mathematical and TeX fidelity}

Natural-language normalization is subordinate to mathematical identity. Protected boundaries include math mode, command names, command options, labels, references, citation keys, paths, comments, and demonstrated original-language arguments. For every Latin/Cyrillic pair, removing Unicode letters and combining marks yields the same non-letter stream. This check preserves braces, control words, operators, numbers, alignment marks, and mathematical symbols.

An earlier failure demonstrates why a linguistic scanner is insufficient. Projector v5 misread the TeX unit \code{24em} as a visible word and created a Cyrillic unit that is not legal TeX. V6 therefore acquired an explicit grammar for numeric dimensions, which v8 inherits and rechecks. It preserves 447 previously exposed occurrences: 219 \code{em} and 228 \code{pt}; Cyrillic dimension-unit residue is zero. The failure remains in the evidence because a clean lexical result does not prove that the source program is executable.

\section{Deterministic v8 derivation}

V8 is a pure function of the four frozen Latin files, the rebased role manifest, and pinned dependencies. It:

\begin{enumerate}
\item authenticates every Latin unit by byte count and SHA-256;
\item authenticates the Latin freeze receipt and freeze decision before projection;
\item verifies exact source and manifest topology;
\item parses TeX into visible and protected roles;
\item merges reviewed boundaries for names, quotations, 722 Roman labels, and dimensions;
\item projects every remaining visible word with the frozen Interslavic transliterator;
\item proves every one of the 372 inherited effects and every one of the 722 Roman loci;
\item licenses every output difference to an exact v020--v021 Latin alignment block;
\item writes output transactionally only after every gate passes, while receipt creation is a separate transaction.
\end{enumerate}

Hostile tests cover clean installation, idempotent replay, third-state refusal, rollback after an injected second-install failure, and refusal of a combined output-and-receipt transaction before projection. Three identical runs pass. An independent verifier that does not import v8 reconstructs all four Cyrillic outputs and then reads the installed bytes back. The test, projector, and verifier are themselves hashed, so ``the same procedure'' means identical code, not merely a similar description.

\section{Build and multilayer validation}

All XeLaTeX operations run serially, with two passes per document, shell escape disabled, and a fixed \code{SOURCE\_DATE\_EPOCH}. Eight editable components produce two complete readers.

\begin{table}[ht]
\centering
\caption{Verified v021 component readers.}\label{tab:en-metrics}
\begin{tabularx}{\textwidth}{@{}YrrY@{}}
\toprule
Artifact & Pages & Bytes & SHA-256\\
\midrule
Latin reader & 565 & 2,993,022 & \hash{57FA897FCADA9C9A8BC764BB0DE30724B7434C31B8C8A35B14FB42922DAB898B}\\
Cyrillic reader & 588 & 3,300,230 & \hash{9513C002794173DD6CC44F157489B0B7A94E92E19378605D3E15C5EE27E895D3}\\
\bottomrule
\end{tabularx}
\end{table}

Release QA reports 38 passing checks and zero errors. It covers exact source readback, TeX and non-letter structure, fatal logs, page topology, fonts, \code{pdftotext}, script census, and text identity between component and cumulative readers. The linked public reader must contain 6 front/divider pages plus 565 Latin and 588 Cyrillic pages, for a total of 1,159. Its builder rejects any other topology.

A structural test does not prove good appearance. The final release gate therefore also requires serial rendering and visual review of the new front matter, every changed page and adjacent control pages, plus the boundaries between scripts. Text extraction alone cannot reveal overlap, clipping, or an unreadable glyph substitution.

\section{Publication, access, and preservation}

The public Zenodo surface contains four products:

\begin{enumerate}
\item \code{00\_NOETHER\_INTERSLAVIC\_COMPLETE\_LINKED\_READER.pdf};
\item \code{01\_NOETHER\_INTERSLAVIC\_EDITABLE\_SOURCES.zip};
\item \code{02\_NOETHER\_INTERSLAVIC\_EVIDENCE\_AND\_PROVENANCE.zip};
\item \code{03\_NOETHER\_INTERSLAVIC\_SHA256\_MANIFEST.txt}.
\end{enumerate}

The first file is the human reader. The source archive contains the Latin canon, derived Cyrillic sources, projector, tests, builders, metadata, and member manifests. The evidence archive contains the append-only editorial ledger, decision sidecars, complete role manifest, validation receipts, this methods paper and its LaTeX source, QA, retained adverse evidence, and its own member manifest. Operational caches, credentials, raw conversation records, and unrelated workspaces are not public evidence products.

The stable edition concept is \doiurl{10.5281/zenodo.21926382}; the v021 version DOI is \doiurl{10.5281/zenodo.22077468}; the public repository is \repo. The repository introduction and Zenodo metadata place a complete Interslavic presentation before a complete English counterpart. Publication concludes with anonymous readback of every public file and comparison of its bytes and SHA-256 with the locally frozen manifest.

\section{Limitations, human review, and responsible claims}

The edition can make strong claims about an exact input vector, replayable editorial history, complete projector-role review, deterministic Cyrillic derivation, bounded output delta, clean builds, and verified PDF structure. It does not claim:

\begin{itemize}
\item native-speaker certification or consensus among Slavic communities;
\item endorsement by the Interslavic community or an official standard;
\item independent scholarly peer review;
\item a new critical edition of the German originals;
\item that computational validation proves idiomaticity or an optimal lexical choice;
\item that Cyrillic is independent evidence for the translation.
\end{itemize}

A human project owner set scope, custody, access, and publication goals. The available evidence does not justify claiming that a human or native speaker reviewed the entire language corpus line by line. ``Completed audit frontier'' in v021 therefore denotes the completed declared occurrence-level and projection checks performed with machine assistance and verified by deterministic controls. It does not denote human linguistic certification.

These limitations are not reasons to hide the work. Public sources, locators, and hashes turn disagreement into an executable correction: a reviewer can identify a locus, propose a better reading, cite contrary evidence, and produce a monotonic successor without erasing the prior state.

\section{Reusable protocol for other languages and transcriptions}

The minimum transferable protocol is:

\begin{enumerate}
\item define one source authority and one target-language canon;
\item assign an explicit role to every derived script, transcription, or format;
\item freeze source units with byte counts and hashes before editing;
\item record material decisions append-only with rejected alternatives;
\item separate visible language from protected syntax, mathematics, and identities;
\item review candidates occurrence by occurrence and bind classifications to stable locators;
\item validate meaning, language, mathematics, syntax, builds, extraction, fonts, and rendering as distinct gates;
\item publish a human reader, editable sources, evidence, and a top-level hash manifest;
\item expose a machine-readable canon index and deterministic entry points;
\item disclose absent review and remaining uncertainty;
\item issue immutable versions and anonymously verify public bytes.
\end{enumerate}

The central lesson is simple: an AI-produced language edition becomes scholarly usable when its choices are not trapped in a model conversation. The canon, roles, evidence, failures, build commands, and public identities must survive independently of the session that created them.

\section*{Citation, availability, and reuse}

Suggested citation: \emph{Machine-assisted Interslavic editorial project. Od jezyčnogo kanona do povtorimoj publikacije: metodologija mašinno-pomaganoj medžuslovjanskoj redakcije / From a language canon to a reproducible publication: machine-assisted Interslavic editing of Emmy Noether's mathematical works. NOETHER-ISV-v021.} \doiurl{10.5281/zenodo.22077468} (2026).

This methods paper and other project-created editorial material are offered under CC0 to the extent that rights exist. Source works and third-party materials retain their own status. The LaTeX source is included so that the paper can be corrected, extended, or republished with accurate attribution, version identity, and an inspectable evidence lineage.

\end{document}
